% Purpose: Comprehensive analysis of regional target datasets
% (SMINO, 2024 Socchieve sequence) and benchmark training corpora
% (INSTANCE, STEAD, SCEDC, Original).
% Status: Draft; under academic review.
% Source-of-truth: OGS/src/ogsconstants.py,
% doc/MHPC/chapters/2_datasets.tex, OGS catalog records.

\label{chap:datasets}

This chapter details the empirical seismological datasets underpinning this
dissertation. Section~\ref{sec:case_study} presents the target operational
domain: the continuous waveform recordings and analyst-curated reference
catalog acquired across Northeastern Italy spanning \textbf{+20 years} of
observational data (\textbf{January 1st, 2005} to
\textbf{December 31st, 2025}), capturing multi-decade background
microseismicity as well as prominent earthquake sequences such as the spring
2024 Socchieve seismic crisis. Section~\ref{sec:benchmark_datasets} presents an
in-depth comparative review of the four benchmark training datasets utilized to
train and evaluate the deep-learning phase pickers: \ac{INSTANCE} (Italy),
\ac{STEAD} (global), \ac{SCEDC} (Southern California), and the original author
training sets of \ac{PhaseNet} and \ac{EQT}. These collections span
localized regional to global tectonic environments and exhibit diverse
\ac{SNR}, noise textures, and labeling standards, establishing a testbed for
evaluating the domain-adaptation hypothesis.

\section{Target Regional Case Study: Northeastern Italy}
\label{sec:case_study}

The primary observational dataset serves as the target domain for evaluating
the end-to-end performance, accuracy, and operational throughput of the
procedural machine learning pipeline. The dataset comprises continuous
multi-channel seismic recordings acquired across Northeastern Italy and
adjacent border sectors in Austria, Slovenia, Veneto, and Trentino-Alto Adige.

The spatial, instrumental, and seismological parameters characterizing the case
study include:
\begin{itemize}
  \item \textbf{Geographical Bounds and Spatial Polygon}: The study region is
    formally bounded by the \ac{OGS} operational polygon
    (\texttt{OGS\_POLY\_REGION}) defined in the central configuration module
    \texttt{ogsconstants} (see Section~\ref{sec:ogsconstants}). The polygon
    encloses an eight-point convex domain spanning longitudes
    $10.0^\circ\,\mathrm{E}$ to $14.5^\circ\,\mathrm{E}$ and latitudes
    $44.5^\circ\,\mathrm{N}$ to $47.0^\circ\,\mathrm{N}$, capturing the active
    thrust fronts of the Eastern Southern Alps and the transcurrent structures
    of the Western Dinaric junction. The extended bounding box for waveform
    queries covers longitudes $9.5^\circ\text{--}15.0^\circ\,\mathrm{E}$ and
    latitudes $44.3^\circ\text{--}47.5^\circ\,\mathrm{N}$
    (\texttt{OGS\_STUDY\_REGION}).
  \item \textbf{Instrumental Network Aperture}: Waveform data were retrieved
    from \textbf{266 distinct seismic stations} belonging to 18 regional,
    national, and international networks (Figure~\ref{fig:ogs_stations}). The
    sensor network comprises high-sensitivity three-component broadband
    velocimeters \eg{Streckeisen STS-2, Nanometrics Trillium 120/240, Guralp
    CMG-3T/40T} and high-dynamic-range strong-motion accelerometers
    (Kinemetrics Episensor). Over 790 continuous channels were monitored,
    providing dense spatial sampling across the epicentral zone and surrounding
    Alpine valleys.
  \item \textbf{Temporal Window and Observation Horizon}: The target
    observation window spans \textbf{+20 years} of continuous observational
    data, from \textbf{January 1st, 2005} to \textbf{December 31st, 2025}
    (00:00:00 UTC to 23:59:59 UTC). This multi-decade horizon captures
    long-term regional background microseismicity, significant earthquake
    sequences \eg{the March 27, 2024 $M_{\mathrm{L}}~4.6$ Socchieve sequence},
    network evolution, and multi-annual seismicity rate variations across the
    Southeastern Alps and northern Adriatic foreland.
  \item \textbf{The March 27, 2024 Socchieve Earthquake}: The focal seismic
    event anchoring this study occurred on \textbf{March 27, 2024, at
    21:19:37 UTC} (22:19:37 local Italian time). The earthquake nucleated in
    the Carnic Prealps approximately $\qty{5}{\kilo\meter}$ south-southwest of
    Socchieve, Udine province (focal parameters: Latitude
      $46.3583^\circ\,\mathrm{N}$, Longitude $12.8080^\circ\,\mathrm{E}$,
      hypocentral depth $\qty{10.22}{\kilo\meter}$, local magnitude
    $M_{\mathrm{L}}~4.6 \pm 0.3$, moment magnitude $M_{\mathrm{w}}~4.5$). This
    event represents the most energetic tectonic earthquake in Friuli since the
    1976--1980 seismic crisis and generated extensive macroseismic shaking
    (EMS-98 intensity V--VI) throughout Friuli and Veneto, triggering
    rockfalls, masonry cracking, and emergency civil protection deployments.
  \item \textbf{Continuous Ingestion and Temporal Availability}: All continuous
    waveform data were retrieved in standardized miniSEED (MSEED) format
    directly from distributed \ac{FDSN} data centers using the automated
    \texttt{ogsdownloader} utility (see Section~\ref{sec:data_acquisition}).
    Data centers queried include \ac{OGS} internal servers,
    \ac{INGV}, GFZ, IRIS,
    ETH, and \ac{ORFEUS}. Station availability over the observation horizon is
    illustrated in Figure~\ref{fig:ogs_availability}, demonstrating sustained
    telemetry uptime across the core regional network. Instrument transfer
    functions, sensor orientations, gains, and geographic coordinates were
    parsed and validated from StationXML metadata files.
  \item \textbf{Sampling Rate Standardization}: All ingested
    continuous waveforms were resampled from their native
    acquisition frequencies (\qtylist{50;100;200}{\hertz}) to
    a standardized frequency of \textbf{\qty{100}{\hertz}} ($\Delta t =
    \qty{0.01}{\second}$) via polyphase anti-aliasing decimation,
    ensuring exact alignment with the input tensors expected by
    SeisBench deep-learning models.
\end{itemize}

\begin{figure}[htbp]
  \centering
  \includegraphics[width=\textwidth]{OGS/ALL/OGSStations.png}
  \caption{Spatial distribution of the 266 seismological and
    accelerometric stations across Northeastern Italy and adjacent
    border regions analyzed in this study. The reader should notice
    the high station density along the Friuli Prealpine belt and the
    transboundary coverage extending into Austria, Slovenia, and the
    Veneto plain. Station coordinates and metadata verified against
  official \ac{FDSN} StationXML inventories from 18 participating networks.}
  \label{fig:ogs_stations}
\end{figure}

\begin{figure}[htbp]
  \centering
  \MissingFigure{The event-map asset required for this figure is not present in
    the repository. A station-distribution map must not be substituted for an
    event-location map. Provide a versioned event-map artifact and its
    generating command before publication.
  }
  \caption{Placeholder for the event-location map within the \ac{OGS}
    operational polygon over the stated observation interval. The
    intended analysis is the spatial relation between analyst-curated
    reference events and the Socchieve sequence; no such relation is
    asserted while the map artifact is unavailable. Evidence
    limitation: the required versioned event-map file and its run
  provenance are absent from this checkout.}
  \label{fig:ogs_seismicity}
\end{figure}

\begin{figure}[htbp]
  \centering
  \includegraphics[width=\textwidth]{OGS/ALL/OGSAvailability.png}
  \caption{Temporal station availability and daily trace acquisition metrics
    across the regional network. The reader should notice the sustained
    operational continuity across participating networks, with minor transient
    drops resolved via automatic multi-client \ac{FDSN} fallback. Comprehensive
    annual station availability distributions covering the complete +20 years
    observational baseline (January 1st, 2005 to December 31st, 2025) across
    all 21 calendar years are detailed in Supplementary
    Material~\ref{sec:supp_annual_ogs_network}. Data generated from
    continuous miniSEED audit logs parsed by \texttt{ogsutils.waveforms}.
  }
  \label{fig:ogs_availability}
\end{figure}

\subsection{Analyst-Curated \ac{OGS} Reference Catalog}
\label{subsec:ogs_ground_truth}

To establish a baseline reference for evaluating the \ac{ML} pipeline,
we used an analyst-curated subset of the revised \ac{OGS} \acf{CRS} seismic
bulletin catalog \textcite{MagrinEtAl2026RevisedOGSCatalog}. The archive is a
reviewed operational reference, not an observation of every physical
earthquake. In standard operational practice, \ac{OGS} seismic analysts
manually pick phase arrivals, assign onset quality weights (see
Table~\ref{tab:h71_weights}), verify event associations, and compute
event hypocenters using Hypo71 or HypoEllipse algorithms. The source files
used here in the supported \texttt{.hpl}, \texttt{.dat}, \texttt{.txt}, and
\texttt{.pun} formats were parsed and merged by the \texttt{DataCatalog}
implementation in \path{OGS/src/ogsparser.py}; its output is a consolidated
catalog of standardized \texttt{EVENTS} and \texttt{PICKS} tables.

However, comparing \ac{ML} hypocenter locations derived
from modern probabilistic non-linear search algorithms against
legacy Hypo71 locations would introduce systematic biases arising
from differing inversion formulations, crustal velocity
discretizations, and station elevation corrections. Therefore, to
isolate the performance of the \ac{ML} picking and
association components, the \ac{OGS} reference events were independently
relocated using the identical \textbf{1D \acf{NonLinLoc}} velocity model and
inversion parameterization employed in our computational pipeline (Magrin,
personal communication).

The resulting relocated \ac{OGS} reference catalog contains:
\begin{itemize}
  \item \textbf{\placeholder{NUM EVENTS} confirmed seismic events} strictly
    contained within the regional polygon boundary \texttt{OGS\_POLY\_REGION}
    during the +20-year observation window (January 1st, 2005 to December 31st,
    2025).
  \item Local magnitudes ranging continuously from microseismic events
    ($M_{\mathrm{L}} = -0.2$) up to the $M_{\mathrm{L}}~4.6$ Socchieve
    mainshock, with completeness threshold $M_{\mathrm{c}} \approx 1.2$.
  \item Focal depths predominantly confined between $\qty{5.0}{\kilo\meter}$
    and $\qty{14.0}{\kilo\meter}$, reflecting the shallow seismogenic brittle
    layer of the Southern Alps.
  \item \textbf{\placeholder{NUM PICKS} manually verified phase picks},
    comprising \textbf{\placeholder{NUM P PICKS} Primary (P)} wave arrivals and
    \textbf{\placeholder{NUM S PICKS} Secondary (S)} wave arrivals. Each pick
    includes analyst-assigned onset polarity, weight, and timing uncertainty.
\end{itemize}

This high-quality reference catalog provides the baseline for the
bipartite graph matching assessment (\ac{BPGMA}) developed in
Chapter~\ref{chap:pipeline_evaluation}.

\section{Benchmark Machine-Learning Training Datasets}
\label{sec:benchmark_datasets}

The deep-learning phase pickers evaluated in this thesis—PhaseNet
and \ac{EQT}—rely on supervised training over massive
corpora of three-component seismic waveforms labeled with expert
ground-truth P- and S-wave arrival times. To understand how models
generalize across varying tectonic regimes and noise backgrounds, we
selected four benchmark datasets available through the SeisBench
framework \parencite{SeisBench}. Each dataset represents a distinct
combination of scale, geographical setting, instrumentation density,
and labeling curation.

\subsection{INSTANCE: The Italian Seismic Dataset for Machine Learning}
\label{subsec:instance}

The \acf{INSTANCE} dataset is a comprehensive open-access collection of labeled
seismic waveforms recorded across the Italian peninsula between 2005 and 2020
by the National Seismic Network (RSN), operated by \ac{INGV}, and partner
regional networks including \ac{OGS}.

Unlike global machine-learning datasets curated primarily for high \ac{SNR},
\ac{INSTANCE} was specifically designed to preserve the realistic operational
conditions of active regional monitoring networks:
\begin{itemize}
  \item \textbf{Operational Noise Realism}: \ac{INSTANCE} incorporates
    the full spectrum of environmental, cultural, and industrial
    noise characteristic of the Mediterranean and Alpine
    environments. This includes industrial vibrations, heavy road
    and rail traffic, quarry blasts, severe meteorological noise,
    and coastal microseisms.
  \item \textbf{Emergent and Low-Magnitude Arrivals}: The dataset
    encompasses a large proportion of small-magnitude events
    ($M_{\mathrm{L}} < 2.5$) exhibiting emergent onsets, complex
    scattering, and low signal-to-noise ratios.
  \item \textbf{Volume and Balance}: \ac{INSTANCE} comprises nearly
    \textbf{1.3 million three-component waveform traces}
    (\num{1291537}), including \num{1159249} labeled P-wave arrivals,
    \num{713883} labeled S-wave arrivals, across \num{54008} distinct events,
    and a dedicated negative class of \num{132288} pure noise windows
    extracted during periods without seismic activity.
  \item \textbf{Geographical Proximity}: Because \ac{INSTANCE} waveforms
    originate from Italy, the training wavefields reflect crustal attenuation,
    complex fault structures, and station vault installations physically
    similar to those of our Northeastern Italy study area.
\end{itemize}

\begin{figure}[htbp]
  \centering
  \includegraphics[width=0.82\textwidth]{INSTANCE.png}
  \caption{Geographical station distribution and event density of the
    \ac{INSTANCE} benchmark dataset across Italy. Notice the
    extensive coverage along the Apenninic chain and the
    Southern Alpine arc, capturing realistic Mediterranean crustal
    paths and anthropogenic noise conditions. Image taken from
  \textcite{SeisBench}.}
  \label{fig:instance_sample}
\end{figure}

\subsection{STEAD: STanford EArthquake Dataset}
\label{subsec:stead}

The \acf{STEAD} dataset \parencite{STEAD} is a globally distributed
benchmark dataset curated by researchers at Stanford University to
facilitate artificial intelligence research in seismology. Sourced
from seismic stations worldwide, \ac{STEAD} incorporates waveforms from
subduction zones, continental collision zones, volcanic systems, and
oceanic intraplate environments.

Key characteristics of \ac{STEAD} include:
\begin{itemize}
  \item \textbf{High Signal Clarity and Curation}: Waveforms were
    subjected to strict automated and manual quality filters,
    ensuring distinct wavefield onsets and elevated signal-to-noise
    ratios for earthquake signals.
  \item \textbf{Extensive Labeling}: \ac{STEAD} contains over
    \textbf{1.26 million three-component traces} (\num{1265657}), with
    \num{1030231} P-wave picks and \num{1030231} S-wave picks across
    \num{441705} distinct earthquakes.
  \item \textbf{Dedicated Noise Partition}: The dataset includes
    \num{235426} continuous noise recordings sampled across global
    stations, allowing neural networks to learn robust
    discriminative boundaries between non-stationary environmental
    noise and true seismic arrivals.
  \item \textbf{Global Transferability}: By encompassing a wide
    variety of sensor types, crustal velocity structures, and
    epicentral distances (local to regional, up to
    $\qty{350}{\kilo\meter}$), \ac{STEAD} is engineered to prevent models from
    overfitting to region-specific propagation paths.
\end{itemize}

\begin{figure}[htbp]
  \centering
  \includegraphics[width=0.82\textwidth]{STEAD.png}
  \caption{Worldwide distribution of seismological stations contributing
    waveform recordings to the global \ac{STEAD} benchmark dataset. The reader
    should notice the broad geographic and tectonic representation across
    diverse orogenic belts and plate boundaries. Image taken from
    \textcite{SeisBench}.
  }
  \label{fig:stead_sample}
\end{figure}

\subsection{SCEDC: Southern California Earthquake Data Center}
\label{subsec:scedc}

The \acf{SCEDC} dataset \parencite{SCEDC} provides an unprecedented
volume of seismic recordings curated from the \acf{SCSN}, spanning decades
of continuous monitoring across the San Andreas fault system.

\ac{SCEDC} is characterized by:
\begin{itemize}
  \item \textbf{Massive Data Volume}: Sourced through SeisBench, the
    \ac{SCEDC} dataset contains \textbf{8.11 million three-component
    traces} (\num{8111060}), incorporating \num{7571970} P-wave picks
    and \num{4364155} S-wave picks across \num{378528} local earthquakes.
  \item \textbf{Dense Local Network Geometry}: The \ac{SCSN} features an
    exceptionally high spatial sensor density, capturing
    high-frequency wavefield details over purely local propagation
    paths ($\Delta \le \qty{150}{\kilo\meter}$).
  \item \textbf{Catalog Completeness}: Decades of continuous analyst
    curation and automated cross-correlation template matching make
    the \ac{SCEDC} catalog complete to very low magnitude thresholds
    ($M_{\mathrm{c}} \approx 1.0$), minimizing missing pick
    annotations in the training labels.
  \item \textbf{Regional Bias}: While mathematically advantageous
    for training deep high-capacity networks with millions of
    parameters without overfitting, the dataset is geographically
    homogeneous, reflecting Southern California crustal velocity
    gradients and strike-slip tectonic structures.
\end{itemize}

\begin{figure}[htbp]
  \centering
  \includegraphics[width=0.82\textwidth]{SCEDC.png}
  \caption{Spatial station distribution of the \ac{SCSN} represented
    in the \ac{SCEDC} dataset. The
    reader should notice the dense spatial grid concentrated along the
    San Andreas fault system and the Los Angeles basin. Image
  taken from \textcite{SeisBench}.}
  \label{fig:scedc_sample}
\end{figure}

\subsection{Original Author Dataset: PhaseNet and EQTransformer}
\label{subsec:original_corpora}

In addition to the three standard community benchmark datasets, we
evaluate the models trained on the original datasets established by
the respective model architects:
\begin{itemize}
  \item \textbf{Original \ac{PhaseNet} Dataset} \parencite{PhaseNet}: Curated
    by the \ac{NCSN}, this dataset
    comprises approximately $779,514$ waveform traces recorded
    across $234,117$ local seismic events in Northern California.
    The dataset features high-quality manual analyst arrival times
    but lacks explicitly segregated noise-only windows, relying
    instead on pre-onset waveform segments to represent ambient noise.
  \item \textbf{Original \ac{EQT} Dataset}
    \parencite{EQTransformer}: Curated as a specialized subset of the
    \ac{STEAD} repository, comprising approximately $1.2$ million
    three-component traces selected to optimize multi-task
    transformer training for simultaneous earthquake detection, P
    picking, and S picking across \qty{60}{\second} waveform windows.
\end{itemize}

\section{Synthesis of Benchmark Datasets and Methodological Comparison}
\label{sec:dataset_synthesis}

Table~\ref{tab:datasets} provides a rigorous, publication-grade
comparative synthesis of all benchmark datasets utilized in this
thesis. All datasets share a common sampling frequency of
$\qty{100}{\hertz}$, either natively or through standardized
SeisBench decimation, ensuring absolute tensor dimensionality
consistency across all deep-learning inference workflows.

\begin{table}[htbp]
  \centering
  \small
  \caption{Systematic comparison of seismological benchmark training
    datasets used for PhaseNet and EQTransformer models. The reader
    should notice the substantial contrast between the massive sample
    volume of SCEDC and the balanced regional noise realism of
    \ac{INSTANCE}. Traces, noise windows, and pick counts are derived from
    SeisBench metadata schemas; distance bounds classify local
    ($\Delta \le \qty{150}{\kilo\meter}$) versus regional ($150 < \Delta \le
  \qty{600}{\kilo\meter}$) propagation paths.}
  \label{tab:datasets}
  \begin{threeparttable}
    \begin{tabularx}{\linewidth}{l r r r r c l}
      \toprule
      \textbf{Dataset} & \textbf{Traces} & \textbf{P Picks} &
      \textbf{S Picks} & \textbf{Events} & \textbf{Noise Traces} &
      \textbf{Coverage / Dist.\tnote{a}} \\
      \midrule
      \ac{INSTANCE} & \num{1291537} & \num{1159249} &
      \num{713883} & \num{54008} & \num{132288} & Italy (L / R) \\
      \ac{STEAD}    & \num{1265657} & \num{1030231} &
      \num{1030231} & \num{441705} & \num{235426} & Global (L / R) \\
      \ac{SCEDC}    & \num{8111060} & \num{7571970} &
      \num{4364155} & \num{378528} & \num{0}\tnote{b} & S. California (L) \\
      \ac{PhaseNet}\tnote{c} & \num{779514} &
      \multicolumn{2}{c}{--- \tnote{d} ---} & \num{234117} &
      \num{0}\tnote{b} & N. California (L) \\
      \ac{EQT}\tnote{c} & \num{1200000} & \num{1000000}
      & \num{1000000} & \num{400000} & \num{200000} & Global (L / R) \\
      \bottomrule
    \end{tabularx}
    \begin{tablenotes}[flushleft]
      \footnotesize
    \item[a] Epicentral distance classification: Local
      ($\mathrm{L}$) denotes $\Delta \le \qty{150}{\kilo\meter}$; Regional
      ($\mathrm{R}$) denotes $\qty{150}{\kilo\meter} < \Delta \le
      \qty{600}{\kilo\meter}$.
    \item[b] Dedicated pure noise trace partitions were not
      independently generated in the \ac{SCEDC} and original \ac{PhaseNet}
      repositories; ambient noise representations are extracted from
      pre-P arrival waveform windows.
    \item[c] Original \ac{PhaseNet} and \ac{EQT} datasets as described in their
      respective publications.
    \item[d] In the original \ac{PhaseNet} publication \parencite{PhaseNet},
      separate counts for P and S arrivals were not reported
      individually; both phase classes are present in the labeled
      training windows.
    \end{tablenotes}
  \end{threeparttable}
\end{table}

\subsection{The Domain Adaptation Hypothesis}
\label{subsec:domain_adaptation_hypothesis}

A central scientific question explored in this dissertation is the
\textbf{domain adaptation hypothesis}: does a deep neural network
trained on massive volumes of data from a distant geological setting
(such as the 8.1 million traces of \ac{SCEDC} in Southern California)
outperform a network trained on a smaller but regionally congruent
dataset (such as the 1.3 million traces of \ac{INSTANCE} in Italy)?

In classical supervised machine learning, statistical theory assumes
that training and testing samples are identically and independently
distributed ($i.i.d.$) from the same underlying probability distribution:
\begin{equation}
  P_{\mathrm{train}}(\mathbf{X}, \mathbf{Y}) =
  P_{\mathrm{test}}(\mathbf{X}, \mathbf{Y}).
  \label{eq:iid_assumption}
\end{equation}
In operational seismology, this assumption is fundamentally violated
when transferring models across geographical regions
($P_{\mathrm{train}} \neq P_{\mathrm{test}}$). Regional differences include:
\begin{enumerate}
  \item \textbf{Path Effects and Velocity Dispersions}: Distinct
    crustal velocity gradients, Conrad/Moho discontinuity depths,
    and intrinsic attenuation ($Q$-factor) modify wavefield
    dispersion and P-to-S coda energy ratios.
  \item \textbf{Anthropogenic and Environmental Noise Spectra}:
    Cultural noise in densely populated European valleys
    (characterized by \qty{50}{\hertz} electrical harmonics, narrow road
    networks, and industrialized river plains) differs substantially
    from noise in the Mojave desert or Southern California coastal basins.
  \item \textbf{Instrumental Hardware Responses}: Variations in
    sensor eigenfrequencies, digitizer bit depths, and vault
    coupling introduce subtle spectral colorations into the
    digitized waveforms.
\end{enumerate}

By systematically evaluating \ac{PhaseNet} and \ac{EQT} models
trained across \ac{INSTANCE}, \ac{STEAD}, \ac{SCEDC}, and the
original datasets on
the identical three-month \ac{SMINO} continuous dataset, this study
provides empirical bounds on regional domain transferability
(presented in Chapter~\ref{chap:results}).
